\documentclass[10pt,a4paper]{amsart}
\usepackage{amsmath,amssymb,amsthm,graphicx}
\usepackage{verbatim}
\newenvironment{reference}{\comment}{\endcomment}
%\newenvironment{reference}{}{}
\newenvironment{slogan}{\comment}{\endcomment}
\newenvironment{history}{\comment}{\endcomment}

% For commutative diagrams we use Xy-pic
\usepackage[all]{xy}

% We use 2cell for 2-commutative diagrams.
\xyoption{2cell}
% Theorem environments.
%
\theoremstyle{plain}
\newtheorem{theorem}[subsection]{Theorem}
\newtheorem{proposition}[subsection]{Proposition}
\newtheorem{lemma}[subsection]{Lemma}

\theoremstyle{definition}
\newtheorem{definition}[subsection]{Definition}
\newtheorem{example}[subsection]{Example}
\newtheorem{exercise}[subsection]{Exercise}
\newtheorem{situation}[subsection]{Situation}

\theoremstyle{remark}
\newtheorem{remark}[subsection]{Remark}
\newtheorem{remarks}[subsection]{Remarks}

\numberwithin{equation}{subsection}

% Macros
%
\def\lim{\mathop{\mathrm{lim}}\nolimits}
\def\colim{\mathop{\mathrm{colim}}\nolimits}
\def\Spec{\mathop{\mathrm{Spec}}}
\def\Hom{\mathop{\mathrm{Hom}}\nolimits}
\def\Ext{\mathop{\mathrm{Ext}}\nolimits}
\def\SheafHom{\mathop{\mathcal{H}\!\mathit{om}}\nolimits}
\def\SheafExt{\mathop{\mathcal{E}\!\mathit{xt}}\nolimits}
\def\Sch{\mathit{Sch}}
\def\Mor{\mathop{\mathrm{Mor}}\nolimits}
\def\Ob{\mathop{\mathrm{Ob}}\nolimits}
\def\Sh{\mathop{\mathit{Sh}}\nolimits}
\def\NL{\mathop{N\!L}\nolimits}
\def\CH{\mathop{\mathrm{CH}}\nolimits}
\def\proetale{{pro\text{-}\acute{e}tale}}
\def\etale{{\acute{e}tale}}
\def\QCoh{\mathit{QCoh}}
\def\Ker{\mathop{\mathrm{Ker}}}
\def\Im{\mathop{\mathrm{Im}}}
\def\Coker{\mathop{\mathrm{Coker}}}
\def\Coim{\mathop{\mathrm{Coim}}}

% Boxtimes
%
\DeclareMathSymbol{\boxtimes}{\mathbin}{AMSa}{"02}

%
% Macros for moduli stacks/spaces
%
\def\QCohstack{\mathcal{QC}\!\mathit{oh}}
\def\Cohstack{\mathcal{C}\!\mathit{oh}}
\def\Spacesstack{\mathcal{S}\!\mathit{paces}}
\def\Quotfunctor{\mathrm{Quot}}
\def\Hilbfunctor{\mathrm{Hilb}}
\def\Curvesstack{\mathcal{C}\!\mathit{urves}}
\def\Polarizedstack{\mathcal{P}\!\mathit{olarized}}
\def\Complexesstack{\mathcal{C}\!\mathit{omplexes}}
% \Pic is the operator that assigns to X its picard group, usage \Pic(X)
% \Picardstack_{X/B} denotes the Picard stack of X over B
% \Picardfunctor_{X/B} denotes the Picard functor of X over B
\def\Pic{\mathop{\mathrm{Pic}}\nolimits}
\def\Picardstack{\mathcal{P}\!\mathit{ic}}
\def\Picardfunctor{\mathrm{Pic}}
\def\Deformationcategory{\mathcal{D}\!\mathit{ef}}
\usepackage{cleveref}
\crefname{originallemma}{original lemma}{original lemmas}
\Crefname{originallemma}{Original lemma}{Original lemmas}
\makeatletter
\newlabel{lemma-Mittag-Leffler}{{03CA}{source}{}{sourcepointer-1}{}}
\newlabel{lemma-Mittag-Leffler@cref}{{[originallemma][0][]03CA}{[1][0][]source}{}{}{}}
\makeatother
\crefname{originallemma}{original lemma}{original lemmas}
\Crefname{originallemma}{Original lemma}{Original lemmas}
\makeatletter
\newlabel{lemma-NAK}{{00DV}{source}{}{sourcepointer-2}{}}
\newlabel{lemma-NAK@cref}{{[originallemma][0][]00DV}{[1][0][]source}{}{}{}}
\makeatother
\crefname{originallemma}{original lemma}{original lemmas}
\Crefname{originallemma}{Original lemma}{Original lemmas}
\makeatletter
\newlabel{lemma-flat-tor-zero}{{00HL}{source}{}{sourcepointer-3}{}}
\newlabel{lemma-flat-tor-zero@cref}{{[originallemma][0][]00HL}{[1][0][]source}{}{}{}}
\makeatother
\title{Four general facts about completion: surjectivity, exactness for flat quotients, and finiteness}
\author{The Stacks Project Authors}
\date{stacks-project:algebra.tex:local-snapshot-20261012; GNU FDL 1.2 or later}
\begin{document}
\maketitle
\noindent\textit{Editorial note.} All delivered mathematical text is the author's own, copied from the pinned Stacks Project chapter file algebra.tex; no statement or proof step was rewritten, reordered or removed.\par\medskip
\noindent\textit{Editorial note.} Preamble provenance, as required for this source: the frozen snapshot stages the chapter's own preamble file beside it (preamble.tex, 217 lines, 8116 bytes, the same shared Stacks-wide preamble the published chow and etale-cohomology cards use); the first line of the pinned chapter file is \textbackslash{}input\{preamble\}, so no preamble text lives inside the chapter itself. The card takes that file's lines 8-18 (the verbatim package and the three \textbackslash{}comment-based environments reference/slogan/history at lines 9-12, section 13.2, together with the Xy-pic setup) and its lines 154 to the runtime end of the file (the native theorem declarations and mathematical macros). Not one character was changed, no package was added, no macro was defined, and the excerpt serves only as preamble, never as a span coordinate.\par\medskip
\noindent\textit{Editorial note.} The same-chapter pointers are located at their \textbackslash{}label lines inside the pinned chapter file; the cross-chapter pointers are located in the project's own tag list staged in the snapshot (sources/stacks-chow/tags.txt), at the verbatim lines carrying their canonical tags.\par\medskip
\noindent\textit{Editorial note.} The counted proof refers to same-chapter statements by label; they are recorded as navigation stubs pointed at their label lines.\par\medskip
\noindent\textit{Editorial note.} The counted spans contain no cite, so the card carries no bibliography.\par\medskip
\noindent\textit{Allowed standard prerequisites (editorial).}\begin{itemize}
\item Nakayama applied over a nilpotent thickening. Invoked in the counted proof at lines 22805-22808. Lifts surjectivity modulo a power of the ideal to surjectivity of the system of quotients.
\item Left exactness of inverse limits. Used in the counted proof (the exactness statement). Gives exactness of the completed sequence except at the last term, which flatness supplies.
\end{itemize}\medskip
\ifdefined\hypertarget\hypertarget{sourcepointer-1}{}\fi
\noindent\textit{Source reference.} Original lemma 03CA: sources/stacks-algebra/pinned-original-source.tex, source lines 21065--21065. Navigation pointer for \textbackslash{}ref\{lemma-Mittag-Leffler\} used inside the counted span (line 22823); the argument that consumes it is retained in the span.\par\medskip
\ifdefined\hypertarget\hypertarget{sourcepointer-2}{}\fi
\noindent\textit{Source reference.} Original lemma 00DV: sources/stacks-algebra/pinned-original-source.tex, source lines 3833--3833. Navigation pointer for \textbackslash{}ref\{lemma-NAK\} used inside the counted span (line 22806); the argument that consumes it is retained in the span.\par\medskip
\ifdefined\hypertarget\hypertarget{sourcepointer-3}{}\fi
\noindent\textit{Source reference.} Original lemma 00HL: sources/stacks-algebra/pinned-original-source.tex, source lines 9254--9254. Navigation pointer for \textbackslash{}ref\{lemma-flat-tor-zero\} used inside the counted span (line 22831); the argument that consumes it is retained in the span.\par\medskip
\begin{lemma}
\label{lemma-completion-generalities}
Let $R$ be a ring. Let $I \subset R$ be an ideal.
Let $\varphi : M \to N$ be a map of $R$-modules.
\begin{enumerate}
\item If $M/IM \to N/IN$ is surjective, then $M^\wedge \to N^\wedge$
is surjective.
\item If $M \to N$ is surjective, then $M^\wedge \to N^\wedge$ is surjective.
\item If $0 \to K \to M \to N \to 0$ is a short exact sequence of
$R$-modules and $N$ is flat, then
$0 \to K^\wedge \to M^\wedge \to N^\wedge \to 0$ is a short exact sequence.
\item The map $M \otimes_R R^\wedge \to M^\wedge$ is
surjective for any finite $R$-module $M$.
\end{enumerate}
\end{lemma}

\begin{proof}
Assume $M/IM \to N/IN$ is surjective. Then the map $M/I^nM \to N/I^nN$
is surjective for each $n \geq 1$ by Nakayama's lemma. More precisely,
apply Lemma \ref{lemma-NAK} part (11) to the
map $M/I^nM \to N/I^nN$ over the ring $R/I^n$ and the nilpotent
ideal $I/I^n$ to see this. Set $K_n = \{x \in M \mid \varphi(x) \in I^nN\}$.
Thus we get short exact sequences
$$
0 \to K_n/I^nM \to M/I^nM \to N/I^nN \to 0
$$
We claim that the canonical map $K_{n + 1}/I^{n + 1}M \to K_n/I^nM$
is surjective. Namely, if $x \in K_n$ write
$\varphi(x) = \sum z_j n_j$ with $z_j \in I^n$, $n_j \in N$.
By assumption we can write $n_j = \varphi(m_j) + \sum z_{jk}n_{jk}$
with $m_j \in M$, $z_{jk} \in I$ and $n_{jk} \in N$. Hence
$$
\varphi(x - \sum z_j m_j) = \sum z_jz_{jk} n_{jk}.
$$
This means that $x' = x - \sum z_j m_j \in K_{n + 1}$ maps
to $x \bmod I^nM$ which proves the claim. Now we may apply
Lemma \ref{lemma-Mittag-Leffler}
to the inverse system of short exact sequences above to see (1).
Part (2) is a special case of (1).
If the assumptions of (3) hold, then for each $n$ the sequence
$$
0 \to K/I^nK \to M/I^nM \to N/I^nN \to 0
$$
is short exact by
Lemma \ref{lemma-flat-tor-zero}.
Hence we can directly apply
Lemma \ref{lemma-Mittag-Leffler}
to conclude (3) is true.
To see (4) choose generators $x_i \in M$, $i = 1, \ldots, n$.
Then the map $R^{\oplus n} \to M$, $(a_1, \ldots, a_n) \mapsto \sum a_ix_i$
is surjective. Hence by (2) we see
$(R^\wedge)^{\oplus n} \to M^\wedge$, $(a_1, \ldots, a_n) \mapsto \sum a_ix_i$
is surjective. Assertion (4) follows from this.
\end{proof}

\end{document}
